% GENERATED FILE. Edit canonical lessons, then run npm run generate:books.
% canonical-source-hash: fnv1a64:8f60023480992c64
% canonical-lessons: TE-C277-phiryadu, TE-C277-mangali, TE-C277-cakali, TE-C277-kummari, TE-C277-kammari, TE-R277-first-pass-words-from-chapters-273-274, TE-R277-second-pass-words-from-chapters-275-277

\chapter{Complaint, Barber, Washerman, Potter, Blacksmith}
\label{ch:te-complaint-barber-washerman-potter-blacksmith}
\hlchaptermodality{277}

\begin{chapteropening}
\textbf{By the end of this chapter:} \emph{I can say a complaint, a barber, a washerman, a potter and a blacksmith.}

Complete the last lesson of chapter 277: I can say a complaint, a barber, a washerman, a potter and a blacksmith.
\end{chapteropening}

\section[\texorpdfstring{\te{ఫిర్యాదు} (phiryādu)}{ఫిర్యాదు (phiryādu)}]{\te{ఫిర్యాదు} (phiryādu) — a complaint}

\label{lesson:TE-C277-phiryadu}



\begin{quote}
Before the new one: say the Telugu for grammar, then the Telugu for a blackboard.
\end{quote}

\subsection*{You'll want to know: \te{ఫిర్యాదు}}
\textbf{\te{ఫిర్యాదు}} — \emph{phiryādu} — \textquotedblleft{}a complaint\textquotedblright{}.

\begin{grammarlens}[title={What you've built}]
One more word for this chapter.
\end{grammarlens}

\subsection*{Your turn}
\begin{itemize}
\raggedright
  \item \emph{Say it:} \emph{phiryādu}
  \item \emph{Say it:} \emph{phiryādu}, once more
  \item \emph{Say it:} say \emph{phiryādu}
  \item \emph{Recall:} say the Telugu for grammar, then the Telugu for a blackboard, then say \emph{phiryādu} again
\end{itemize}

\begin{tcolorbox}[breakable,colback=teal!4,colframe=teal!35!black,title={Before you move on}]
What does \te{ఫిర్యాదు} mean? (\textquotedblleft{}A complaint\textquotedblright{}.) Say it once more.
\end{tcolorbox}
\section[\texorpdfstring{\te{మంగలి} (maṅgali)}{మంగలి (maṅgali)}]{\te{మంగలి} (maṅgali) — a barber}

\label{lesson:TE-C277-mangali}



\begin{quote}
Before the new one: say the Telugu for a blackboard, then the Telugu for a complaint.
\end{quote}

\subsection*{You'll want to know: \te{మంగలి}}
\textbf{\te{మంగలి}} — \emph{maṅgali} — \textquotedblleft{}a barber\textquotedblright{}.

\begin{grammarlens}[title={What you've built}]
One more word for this chapter.
\end{grammarlens}

\subsection*{Your turn}
\begin{itemize}
\raggedright
  \item \emph{Say it:} \emph{maṅgali}
  \item \emph{Say it:} \emph{maṅgali}, once more
  \item \emph{Say it:} say \emph{maṅgali}
  \item \emph{Recall:} say the Telugu for a blackboard, then the Telugu for a complaint, then say \emph{maṅgali} again
\end{itemize}

\begin{tcolorbox}[breakable,colback=teal!4,colframe=teal!35!black,title={Before you move on}]
What does \te{మంగలి} mean? (\textquotedblleft{}A barber\textquotedblright{}.) Say it once more.
\end{tcolorbox}
\section[\texorpdfstring{\te{చాకలి} (cākali)}{చాకలి (cākali)}]{\te{చాకలి} (cākali) — a washerman}

\label{lesson:TE-C277-cakali}



\begin{quote}
Before the new one: say the Telugu for a complaint, then the Telugu for a barber.
\end{quote}

\subsection*{You'll want to know: \te{చాకలి}}
\textbf{\te{చాకలి}} — \emph{cākali} — \textquotedblleft{}a washerman\textquotedblright{}.

\begin{grammarlens}[title={What you've built}]
One more word for this chapter.
\end{grammarlens}

\subsection*{Your turn}
\begin{itemize}
\raggedright
  \item \emph{Say it:} \emph{cākali}
  \item \emph{Say it:} \emph{cākali}, once more
  \item \emph{Say it:} say \emph{cākali}
  \item \emph{Recall:} say the Telugu for a complaint, then the Telugu for a barber, then say \emph{cākali} again
\end{itemize}

\begin{tcolorbox}[breakable,colback=teal!4,colframe=teal!35!black,title={Before you move on}]
What does \te{చాకలి} mean? (\textquotedblleft{}A washerman\textquotedblright{}.) Say it once more.
\end{tcolorbox}
\section[\texorpdfstring{\te{కుమ్మరి} (kummari)}{కుమ్మరి (kummari)}]{\te{కుమ్మరి} (kummari) — a potter}

\label{lesson:TE-C277-kummari}



\begin{quote}
Before the new one: say the Telugu for a barber, then the Telugu for a washerman.
\end{quote}

\subsection*{You'll want to know: \te{కుమ్మరి}}
\textbf{\te{కుమ్మరి}} — \emph{kummari} — \textquotedblleft{}a potter\textquotedblright{}.

\begin{grammarlens}[title={What you've built}]
One more word for this chapter.
\end{grammarlens}

\subsection*{Your turn}
\begin{itemize}
\raggedright
  \item \emph{Say it:} \emph{kummari}
  \item \emph{Say it:} \emph{kummari}, once more
  \item \emph{Say it:} say \emph{kummari}
  \item \emph{Recall:} say the Telugu for a barber, then the Telugu for a washerman, then say \emph{kummari} again
\end{itemize}

\begin{tcolorbox}[breakable,colback=teal!4,colframe=teal!35!black,title={Before you move on}]
What does \te{కుమ్మరి} mean? (\textquotedblleft{}A potter\textquotedblright{}.) Say it once more.
\end{tcolorbox}
\section[\texorpdfstring{\te{కమ్మరి} (kammari)}{కమ్మరి (kammari)}]{\te{కమ్మరి} (kammari) — a blacksmith}

\label{lesson:TE-C277-kammari}



\begin{quote}
Before the new one: say the Telugu for a washerman, then the Telugu for a potter.
\end{quote}

\subsection*{You'll want to know: \te{కమ్మరి}}
\textbf{\te{కమ్మరి}} — \emph{kammari} — \textquotedblleft{}a blacksmith\textquotedblright{}.

\begin{grammarlens}[title={What you've built}]
One more word for this chapter.
\end{grammarlens}

\subsection*{Your turn}
\begin{itemize}
\raggedright
  \item \emph{Say it:} \emph{kammari}
  \item \emph{Say it:} \emph{kammari}, once more
  \item \emph{Say it:} say \emph{kammari}
  \item \emph{Recall:} say the Telugu for a washerman, then the Telugu for a potter, then say \emph{kammari} again
\end{itemize}

\begin{tcolorbox}[breakable,colback=teal!4,colframe=teal!35!black,title={Before you move on}]
What does \te{కమ్మరి} mean? (\textquotedblleft{}A blacksmith\textquotedblright{}.) Say it once more.
\end{tcolorbox}
\section[(dialogue)]{Words from chapters 273-274 — a first pass}

\label{lesson:TE-R277-first-pass-words-from-chapters-273-274}



\begin{quote}
Say the words for a potter and a blacksmith.
\end{quote}

\subsection*{Your turn}
Say these aloud:

\begin{itemize}
\raggedright
  \item interest, a bill, a discount, groceries, a purchase
  \item a class, a lesson, an answer, a subject, mathematics
\end{itemize}

\begin{tcolorbox}[breakable,colback=teal!4,colframe=teal!35!black,title={Before you move on}]
Interest? (\textbf{\te{వడ్డీ}}, \emph{vaḍḍī}.) A bill? (\textbf{\te{బిల్లు}}, \emph{billu}.) A discount? (\textbf{\te{తగ్గింపు}}, \emph{taggimpu}.) Groceries? (\textbf{\te{సరుకులు}}, \emph{sarukulu}.) A purchase? (\textbf{\te{కొనుగోలు}}, \emph{konugōlu}.) A class? (\textbf{\te{తరగతి}}, \emph{taragati}.) A lesson? (\textbf{\te{పాఠం}}, \emph{pāṭhaṁ}.) An answer? (\textbf{\te{జవాబు}}, \emph{javābu}.) A subject? (\textbf{\te{విషయం}}, \emph{viṣayaṁ}.) Mathematics? (\textbf{\te{గణితం}}, \emph{gaṇitaṁ}.)
\end{tcolorbox}
\section[(dialogue)]{Words from chapters 275-277 — a second pass}

\label{lesson:TE-R277-second-pass-words-from-chapters-275-277}



\begin{quote}
Say the words for a priest and an artist.
\end{quote}

\subsection*{Your turn}
Say these aloud:

\begin{itemize}
\raggedright
  \item a priest, an artist, a poet, music, a play
  \item chemistry, biology, economics, grammar, a blackboard
  \item a complaint, a barber, a washerman, a potter, a blacksmith
\end{itemize}

\begin{tcolorbox}[breakable,colback=teal!4,colframe=teal!35!black,title={Before you move on}]
A priest? (\textbf{\te{పూజారి}}, \emph{pūjāri}.) An artist? (\textbf{\te{కళాకారుడు}}, \emph{kaḷākāruḍu}.) A poet? (\textbf{\te{కవి}}, \emph{kavi}.) Music? (\textbf{\te{సంగీతం}}, \emph{saṅgītaṁ}.) A play? (\textbf{\te{నాటకం}}, \emph{nāṭakaṁ}.) Chemistry? (\textbf{\te{రసాయనశాస్త్రం}}, \emph{rasāyanaśāstraṁ}.) Biology? (\textbf{\te{జీవశాస్త్రం}}, \emph{jīvaśāstraṁ}.) Economics? (\textbf{\te{అర్థశాస్త్రం}}, \emph{arthaśāstraṁ}.) Grammar? (\textbf{\te{వ్యాకరణం}}, \emph{vyākaraṇaṁ}.) A blackboard? (\textbf{\te{నల్లబల్ల}}, \emph{nallaballa}.) A complaint? (\textbf{\te{ఫిర్యాదు}}, \emph{phiryādu}.) A barber? (\textbf{\te{మంగలి}}, \emph{maṅgali}.) A washerman? (\textbf{\te{చాకలి}}, \emph{cākali}.) A potter? (\textbf{\te{కుమ్మరి}}, \emph{kummari}.) A blacksmith? (\textbf{\te{కమ్మరి}}, \emph{kammari}.)
\end{tcolorbox}
